% chapters/ch10.tex —— 第 10 章
% 体例：具体失败 → 设计问题 → 机制解释 → 对照实验 → 条件变化 → 选择原则
% 数字身份纪律：
%   - 2506.10943（SEAL）、2506.01716、2505.03335、2607.13104 的数字与判断一律 \cited{}；
%   - 本书 E01/E03c 的读数一律 \measured{}；
%   - 由上述数字相减得到的差值一律 \calcnum{}。
% E08（经验留在外部还是写入参数）未运行，本章不给参数方案的任何实测数字。
\chapter{从外部经验到模型参数}
\label{ch10}

\chaptergoal{什么时候值得把经验写进 $\theta$，而不只是写进 $M$、$P$ 或 $L$？}
\chapterprops{本章讨论写入位置如何改变迁移、归因与恢复条件。P3b关心拟合目标与跨读取者迁移，P6关心更新退化与幅度；它们不能简化成参数位置的一条必然规律。
E08尚未运行，以下以来源结果、机制比较与教学计算展开。}
\chapterfig{图~\ref{fig:ch10-where}：经验的四个写入位置、各自的代价与可逆性}
\chapterpitfall{一段接一段地把新知识微调进模型，每一段单独看都涨了，
再问一周前那一段的题目时它答不上来}

% ---------- 1. 具体失败 ----------
\begin{failurebox}
SEAL 在知识注入设置中先根据文档生成自我编辑文本，再用这些文本和原文进行参数适配；在另一项少样本推理设置中，自我编辑还可以指定数据增强与训练配置。两项设置不能拼成同一个已运行流程\cite{seal-2506}。知识注入的单段 LoRA 实验中，Qwen2.5-7B 在无原文问答上的成绩从只训练原文的\cited{33.5\%}提高到 SEAL 的\cited{47.0\%}。

学进了新材料，先前学到的内容还能留下多少？SEAL 在每条序列中依次吸收八段材料，每次编辑后重测已读的段落。Figure~6 把不同序列中排在相同位置的材料放在同一列。首位材料的平均成绩在第一次编辑后为\cited{0.38}，到第八次编辑后降为\cited{0.19}\cite{seal-2506}。要看懂这一变化，既要选对参照，也要知道每个数平均了哪些材料。

作者说明当前训练未专门优化保持能力，同时提出回归奖励、持续学习约束等后续方向。没有在本项实验中证明这些补救有效，不等于没有讨论补救。这个案例给本章的起点是：一次适配的收益与多次适配后的旧任务表现，需要在同一条版本序列上分别观察。
\end{failurebox}

\designquestion{外部条目、提示、技能与参数具有不同的修改和恢复方式。外部对象通常更容易逐项查看，但删除一个源条目未必撤销派生摘要或已经执行的动作；参数也可以通过保留检查点或独立适配器恢复。应比较具体的更新、使用和撤销条件，而非预先规定一种位置永远更好。}

本章先比较写入位置，再以 SEAL 的两层学习过程说明经验怎样成为训练材料，随后讨论反馈、回归与部署成本。E08 尚无可训练模型、硬件和完整冻结协议；当前实验台未实现参数训练。本章不报告本书自己的参数方案实测。

% ---------- 2. 四个写入位置 ----------
\section{经验的四个写入位置}
\label{sec:ch10-where}

\begin{sloppypar}
本节聚焦 $M$（经验记忆）、$P$（指令与策略）、$L$（工具与技能）与 $\theta$（参数）；工作流 $G$ 的修改另见第~\ref{ch09}章。四个位置不是必经阶梯，费用也不保证依次增加。综述将脚手架更新概括为通常较快、更易逆转，将参数巩固视为可能丢失细节的压缩，并强调更新后的重新验证\cite{self-improvements-survey-2607}。本书把这些观点作为设计动机，不将其改写成每次参数更新必然遗忘，或外部更新必然保持安全的定律。\end{sloppypar}

% figures/ch10-where.tex —— 第 10 章核心图（全书引用入口）
% 用法：在 book/chapters/ch10.tex 中 % figures/ch10-where.tex —— 第 10 章核心图（全书引用入口）
% 用法：在 book/chapters/ch10.tex 中 % figures/ch10-where.tex —— 第 10 章核心图（全书引用入口）
% 用法：在 book/chapters/ch10.tex 中 \input{figures/ch10-where}
% 图主体在 ch10-where-body.tex，便于单独编译与复用。
\begin{figure}[htbp]
\centering
\input{figures/ch10-where-body}
\caption{四种写入位置的记录与恢复条件。它们不是费用必然递增的阶梯；版本可恢复也不等于外部影响已撤销。}
\label{fig:ch10-where}
\end{figure}

% 图主体在 ch10-where-body.tex，便于单独编译与复用。
\begin{figure}[htbp]
\centering
% Four update locations; no assumed monotonic cost or safety ranking.
\begin{tikzpicture}[
 box/.style={draw,rounded corners=2pt,text width=28mm,minimum height=14mm,align=center,font=\small},
 note/.style={text width=28mm,align=left,font=\footnotesize}]
 \node[box,fill=blue!8] at (0,0) {$M$\\经验记忆};
 \node[box,fill=blue!8] at (3.4,0) {$P$\\指令策略};
 \node[box,fill=orange!10] at (6.8,0) {$L$\\工具技能};
 \node[box,fill=orange!10] at (10.2,0) {$\theta$\\模型参数};
 \node[note,anchor=north] at (0,-1) {保存条目与来源\\核对派生摘要\\撤销后检查读取};
 \node[note,anchor=north] at (3.4,-1) {保存提示版本\\核对工具权限\\回归实际行为};
 \node[note,anchor=north] at (6.8,-1) {保存源码与依赖\\验收调用契约\\记录维护成本};
 \node[note,anchor=north] at (10.2,-1) {保存基座与更新\\检查点或适配器\\核对训练及回归};
 \draw[gray!50] (-1.4,-3) -- (11.6,-3);
 \node[align=left,font=\footnotesize,text width=128mm,anchor=north] at (5.1,-3.2)
 {共同问题：更新了什么、依据来自哪里、代价覆盖哪些阶段、恢复能撤销什么。\\
 外部对象通常更易定位；参数可保存版本。两者都需要行为与风险检查，费用不按位置预定排序。};
\end{tikzpicture}

\caption{四种写入位置的记录与恢复条件。它们不是费用必然递增的阶梯；版本可恢复也不等于外部影响已撤销。}
\label{fig:ch10-where}
\end{figure}

% 图主体在 ch10-where-body.tex，便于单独编译与复用。
\begin{figure}[htbp]
\centering
% Four update locations; no assumed monotonic cost or safety ranking.
\begin{tikzpicture}[
 box/.style={draw,rounded corners=2pt,text width=28mm,minimum height=14mm,align=center,font=\small},
 note/.style={text width=28mm,align=left,font=\footnotesize}]
 \node[box,fill=blue!8] at (0,0) {$M$\\经验记忆};
 \node[box,fill=blue!8] at (3.4,0) {$P$\\指令策略};
 \node[box,fill=orange!10] at (6.8,0) {$L$\\工具技能};
 \node[box,fill=orange!10] at (10.2,0) {$\theta$\\模型参数};
 \node[note,anchor=north] at (0,-1) {保存条目与来源\\核对派生摘要\\撤销后检查读取};
 \node[note,anchor=north] at (3.4,-1) {保存提示版本\\核对工具权限\\回归实际行为};
 \node[note,anchor=north] at (6.8,-1) {保存源码与依赖\\验收调用契约\\记录维护成本};
 \node[note,anchor=north] at (10.2,-1) {保存基座与更新\\检查点或适配器\\核对训练及回归};
 \draw[gray!50] (-1.4,-3) -- (11.6,-3);
 \node[align=left,font=\footnotesize,text width=128mm,anchor=north] at (5.1,-3.2)
 {共同问题：更新了什么、依据来自哪里、代价覆盖哪些阶段、恢复能撤销什么。\\
 外部对象通常更易定位；参数可保存版本。两者都需要行为与风险检查，费用不按位置预定排序。};
\end{tikzpicture}

\caption{四种写入位置的记录与恢复条件。它们不是费用必然递增的阶梯；版本可恢复也不等于外部影响已撤销。}
\label{fig:ch10-where}
\end{figure}


\begin{sloppypar}
同一篇还给了一对本书采用的术语来源：把一个系统的自我修改分成\textbf{分布层自指}
（自生成的经验经外部优化写回 $\theta$）与\textbf{动作层自指}
（系统直接编辑 $p,m,\mathcal T,g$，即本书的 $P,M,L,G$）（\cite{self-improvements-survey-2607} §3.2 式 (5)(6)）。
前九章主要讨论外部持久对象的更新，本章进一步讨论自生成材料参与参数学习；具体系统可以同时使用两种路径。
\end{sloppypar}

\begin{table}[htbp]
\centering
\caption{四个写入位置：写入代价、能否逐条追责、换读取者后的表现}
\label{tab:ch10-where}
\begin{small}
\setlength{\tabcolsep}{3pt}
\begin{tabular}{>{\raggedright\arraybackslash}p{1.6cm}>{\raggedright\arraybackslash}p{3.6cm}>{\raggedright\arraybackslash}p{2.4cm}>{\raggedright\arraybackslash}p{4.0cm}>{\raggedright\arraybackslash}p{2.6cm}}
\toprule
位置 & 写一次要付出什么 & 能定位修改对象吗 & 换读取者（换模型 / 换 harness）后 & 本书的实测状态 \\
\midrule
$M$ 经验记忆 & 一条条准入、再检索 & 可定位条目；因果需另测 & 需区分拟合目标与读取者；本书 E03c 观察到 \measured{1}/\measured{16}（求解器拟合）与 \measured{2}/\measured{16}（问题域拟合）道题的收益点估计由正转负 & E03与E03c共十臂有记录；拆分A的采用范围见第\ref{ch05}章\S\ref{sec:ch05-e03} \\
\midrule
$P$ 指令策略 & 改动进入下一次推理 & 可定位版本；因果需另测 & 本书未见同协议对照 & 未跑对照（E04 标受限） \\
\midrule
$L$ 工具技能 & 构建、验收与维护 & 可定位技能；因果需另测 & 本书未见同协议对照 & E05仅工程试跑，无判决 \\
\midrule
$\theta$ 参数 & 执行优化与范围明确的回归 & 难逐条对应；可存检查点或适配器 & 需区分训练材料迁移、参数迁移与换任务 & 未跑（E08 待模型与硬件） \\
\bottomrule
\end{tabular}
\end{small}
\end{table}

\begin{sloppypar}
表中反号计数由每题五次运行得到，仅描述估计符号，没有检验逐题真实效应的符号。表中状态说明本书现有证据覆盖哪里，不提供四个位置的效率排名。参数变化依附于指定基座和架构，不能像一段文本那样直接交给任意读取者；但由经验生成的训练材料仍可用于其他模型，适配器也可在兼容条件下研究迁移。换任务、换接口、换基座分别改变不同条件，不存在一个可以替代所有迁移问题的“唯一等价检验”。\end{sloppypar}

\begin{sloppypar}
还有一条实测读数必须放在这里：本书 E03c 把静态基线从同模型换到换模型，
封存集通过率从 \measured{0.6875} 掉到 \measured{0.2625}。
它说明在这个固定任务范围内，两种读取者的静态基线不同；分析经验迁移增量应分别减去相应基线，不能把差异泛化为模型在所有任务上变弱。
这条纪律在 $\theta$ 方案里同样成立，只是把\textbf{读取者}换成了\textbf{那个被反复更新后的模型}。
\end{sloppypar}

% ---------- 3. 轨迹蒸馏与监督微调 ----------
\subsection{事实更新、知识编辑与能力学习的边界}
\label{sec:ch10-boundary}

\begin{sloppypar}
\textbf{写进 $\theta$}这句话容易把两件事混成一件：一是\textbf{事实更新}（知道一件新事），
二是\textbf{能力学习}（会做一类新事）。两者的成功标准不同，失效方式也不同。
\end{sloppypar}

\begin{table}[htbp]
\centering
\caption{两类写入：问什么、指望什么、失败长什么样}
\label{tab:ch10-kind}
\begin{small}
\setlength{\tabcolsep}{3pt}
\begin{tabular}{>{\raggedright\arraybackslash}p{2.2cm}>{\raggedright\arraybackslash}p{3.8cm}>{\raggedright\arraybackslash}p{4.4cm}>{\raggedright\arraybackslash}p{4.2cm}}
\toprule
类别 & 典型目标 & 成功长什么样 & 失败长什么样 \\
\midrule
事实更新 & 问起上周那条新信息 & 不给原文也答得出 & 新旧条件混淆，或旧任务保持下降；仅凭成绩不能断言具体事实被覆盖 \\
\midrule
能力学习 & 给几个示范就会做新任务 & 少样本成功率上升 & 只记住示范的表面形式，换题型就失效 \\
\bottomrule
\end{tabular}
\end{small}
\end{table}

\begin{sloppypar}
前面那组实验正好两侧各占一半：知识注入那一侧收的是事实
（无原文问答），少样本那一侧评价的是筛选范围内新任务的适配；这只是能力的一种操作化，不穷尽所有能力学习。
\textbf{两侧各自都留下了没解决的部分}：事实侧把 7B 模型从 \cited{32.7} 提到 \cited{47.0}，
但那是\textbf{每段一次更新}这一档的成绩；另一个连续更新实验显示旧任务保持风险，不能从两种设置推出每次连续吸收必然退化。
能力侧还需检验未筛任务与其他基座，不能从这一小范围结果推出普遍适配能力。
\end{sloppypar}

\begin{sloppypar}
这条边界对工程的含义很直接：\textbf{连续参数更新需要预先安排版本、回归与恢复条件}。
若事实频繁变化且需要来源追溯，外部存储通常值得优先比较；重复能力是否应固化为技能或参数，则取决于表达、资源与评价条件，不能规定所有系统必须按同一顺序升级。
\end{sloppypar}

% ---------- 6. 三道验证 ----------
\section{轨迹蒸馏与监督微调}
\label{sec:ch10-sft}

\begin{sloppypar}
把经验写进 $\theta$ 最简单的一条路是：先把经验变成一张训练表，再做一次监督微调。
还需确定训练材料的生成与筛选、更新参数范围、损失权重、优化预算和回归评价；其中任何一项都可能改变结果。
\end{sloppypar}

\begin{sloppypar}
SEAL 的外层学习优化自我编辑生成策略，内层根据编辑材料或配置进行适配，再把适配后的任务表现用于外层选择。知识注入内层采用 LoRA 的设置与多文档全参数训练需要分开。作者报告一次自我编辑评价约\cited{30--45}秒，包含适配与评价。本节据此讨论知识注入的预算，硬件与逐轮计算条件见 Appendix~B.5。少样本 TTT 的时长与硬件另见 Appendix~A.4\cite{seal-2506}。它不能与本书一个未对齐的经验生成时长相除，就宣布两种路径差两个数量级。\end{sloppypar}

\begin{table}[htbp]
\centering
\caption{自造训练表在两个场景上的成绩（均为原论文报告值）}
\label{tab:ch10-seal}
\begin{small}
\begin{tabular}{>{\raggedright\arraybackslash}p{3.6cm}>{\raggedright\arraybackslash}p{4.6cm}>{\raggedright\arraybackslash}p{5.6cm}}
\toprule
场景 & 设置 & 结果 \\
\midrule
少样本抽象推理 & Llama-3.2-1B-Instruct；训练期每题采样 \cited{15} 个候选自我编辑，评测期每题 \cited{5} 个；按编辑成功比例评价 & 上下文学习 \cited{0}、测试时训练 + 自我编辑（未经 RL）\cited{20}、SEAL \cited{72.5}、筛选条件下 Oracle TTT \cited{100} \\
\midrule
知识注入（单段） & Qwen2.5-7B；LoRA；无原文对照问答 & 基座 \cited{32.7}、只在原文上微调 \cited{33.5}、加自生成数据（基座产出）\cited{39.7}、加 GPT-4.1 生成数据 \cited{46.3}、SEAL \cited{47.0} \\
\bottomrule
\end{tabular}
\end{small}
\end{table}

\begin{sloppypar}
表中两种场景要连着读。第一行说明\textbf{路子怎么走}本身是可学的：
在预先筛选的可被Oracle TTT解决的题目中，有没有学过\textbf{改哪些数据、用什么超参数}，
编辑评价成功比例是 \cited{20} 与 \cited{72.5} 的差别。第二行说明另一面：
\textbf{这项单段设置中，原文训练增量很小}（\cited{32.7} 到 \cited{33.5}），
增加合成材料后，点估计更高。知识注入问答由固定版本 GPT-4.1 按参考答案作 yes/no 判定，模型标识与判题提示见原文 Appendix~B.4\cite{seal-2506}；
而把\textbf{怎么重组}也交给学习，又能再往上走一档
（\cited{39.7} 到 \cited{47.0}，差 \calcnum{+7.3} 个百分点）。
\end{sloppypar}

\begin{sloppypar}
原文还报告了多文档继续预训练两列，训练方式与单段不同；
它们是决定\textbf{要不要写 $\theta$}时必须看的数字：
\end{sloppypar}

\begin{table}[htbp]
\centering
\caption{不同训练设置下的点估计：单段与多段不能当纯规模消融}
\label{tab:ch10-cpt}
\begin{tabular}{>{\raggedright\arraybackslash}p{5.6cm}cc}
\toprule
方法 & 全参数训练 $n=200$ & 全参数训练 $n=2067$ \\
\midrule
只在原文上训练 & \cited{36.0} & \cited{31.2} \\
加自生成数据（SEAL） & \cited{58.2} & \cited{46.4} \\
加 GPT-4.1 生成数据 & \cited{59.4} & \cited{49.2} \\
两者之差（本书按原表两行相减） & \calcnum{$-1.2$} & \calcnum{$-2.8$} \\
\bottomrule
\end{tabular}
\end{table}

\begin{sloppypar}
单段设置中 SEAL 相对 GPT-4.1 合成材料高\calcnum{0.7}个百分点，多段两档分别低\calcnum{1.2}与\calcnum{2.8}个百分点。原表单段使用 LoRA，多段使用全参数训练，较大档的基座成绩和评价范围也不同\cite{seal-2506}。因此这三个差值不是只改变语料数量的受控剂量曲线，不能据此断言自生成材料的收益必随规模单调下降。自生成也需要采样与筛选，并非免除数据生成成本。\end{sloppypar}

\begin{sloppypar}
Self-Challenging Language Model Agents 将任务构造成指令、验证函数、示例解与失败案例，用可执行检查筛掉部分不可行或含糊任务，再训练执行者。原文 §4.3 在同策略采样、二值结果奖励的一步 REINFORCE 设置下，将正奖励轨迹的对数似然更新与拒绝采样微调相联系\cite{self-challenging-2506}。该联系需要注明样本分布、归一化与更新项；它不是任意 RL 算法都与 SFT 等价，更不能据奖励只有零和一就省去算法比较。\end{sloppypar}

\subsection{SEAL 的两层学习分别更新了什么}
\label{sec:ch10-sealloops}

读 SEAL 时，可以先把一次自我编辑展开成三个对象：产生编辑的模型版本、编辑所指定的训练材料或配置，以及适配后的候选模型。外层希望改善第一个对象的生成策略，内层则把第二个对象用于第三个对象的参数更新。适配后回答某道题成功，既不意味着编辑文本中的每句话都正确，也不意味着外层策略已经能在所有新材料上生成有效编辑。两层分别需要评价。

外层梯度还采用了一项近似。适配评分随编辑内容变化，也直接依赖作为适配起点的当前模型参数。SEAL 求外层梯度时把本轮评分视为固定量，用它加权编辑文本的对数似然梯度，原文式(3)因此给出近似梯度，式(4)再将其展开为逐token的求和\cite{seal-2506}。内层参数适配与外层更新在这里通过评分连接，适配过程的导数没有沿整条路径传回外层。附录~\ref{sec:appc-pg} 给出的固定奖励定义下的策略梯度恒等式提供了辨认这项近似的参照。

在知识注入设置中，模型接收一段文档，生成蕴含、复述等合成文本；内层把这些文本与原文用于训练，问答时不再把原文放入上下文。因而无原文问答上升，支持的是这套生成与适配流程在指定评价中的效果。它没有单独证明生成的每一条蕴含都由原文严格推出，也没有把收益分解成新增事实、重复曝光、措辞变化或训练长度各自的贡献。要作这样的归因，需要额外控制。

少样本抽象推理中的编辑对象不同：它指定允许的数据增强与训练配置，包括学习率、训练轮次和损失计算范围\cite{seal-2506}。这个编辑更接近一个受限的训练方案。方案中出现某个学习率，不代表模型已经掌握了关于最优学习率的普遍知识；外层只观察它在当前任务条件下的表现。迁移到更大模型、不同数据量或其他优化器时，这份配置可能失效。

原文的抽象推理任务也经过了明确筛选。作者从 ARC 训练与评价范围中分别取\cited{11}与\cited{8}个任务，要求这些任务能够被所用基座在 Oracle TTT 配置下解决；评价时每题生成\cited{5}份编辑，分别适配，统计产生正确输出的编辑比例\cite{seal-2506}。所以\cited{72.5\%}的分母是这份筛选范围内的编辑试验，不是原始 ARC 全集的任务解决率，也不是从五份中免费挑出最好一份后的成功率。

按\cited{8}题、每题\cited{5}份编辑计算，共有\calcnum{40}次编辑评价，\cited{72.5\%}对应\calcnum{29}次成功。这一算术帮助读者理解分母，却不能把四十次都视为独立任务族。相同任务上的编辑共享题目结构，区间与迁移推断应考虑这种聚集。Oracle 的\cited{100\%}也包含预先筛选的条件，不是人类配置在所有未筛任务上必然满分。

这种设计可以检验在可适配任务上学会配置的程度。若要进一步回答“它能发现原本不可适配任务的训练方案吗”，则需要保留相应任务范围；若要回答“一个新任务只允许一次编辑时表现如何”，就要对应一次生成的评价。研究问题不同，并不要求否定原实验；清楚的分母反而使已有结果更容易被正确使用。

外层选择还会消耗比最终部署更大的预算。SEAL 的知识注入 ReSTEM 轮次中，作者对\cited{50}段材料各生成\cited{5}份编辑，每份用\cited{3}个随机种子评价并取均值，再为每段材料保留均值最高的一份编辑，合计\calcnum{750}次内层过程，报告一轮约在\cited{2}张 H100 上耗时\cited{6}小时\cite{seal-2506}。这些是原论文的特定实现成本，不是本书硬件上的预测。只统计最终保留的五十份编辑，会漏掉选择它们所付出的主要代价。

外层学会生成编辑后，部署一个新文档时仍需要内层适配。不能因为外层训练已经完成，就把每份新材料的优化过程当成免费推理；也不能反过来把一次外层训练全额加到每次部署上。合理做法是分别记录学习编辑策略的一次性投入、新材料适配费，以及之后的服务和回归费，再在明确复用范围内讨论摊销。

\subsection{一张带数字的热图怎样描述遗忘}
\label{sec:ch10-retentionmatrix}

SEAL Figure~6 的行表示已经完成的编辑数，列表示材料在序列中的位置。作者公布了\cited{8}组运行摘要，每组汇总\cited{8}条序列，每条又从评价材料中抽取\cited{8}段\cite{seal-continual-archive-2025}。第一列汇总的是各条序列中排在首位的材料。沿这一列向下，比较同一批首位材料在不同编辑步的平均成绩。沿同一行向右，材料的位置和距其首次学习的步数都在变。零次编辑时先测全部材料，此后每次只回测已学位置；上三角尚未回测的位置应留空，归档数组在这些位置填入的零不参与成绩比较。

以第一列为例，初始读数为\cited{0.33}，第一次编辑后为\cited{0.38}，第八次编辑后为\cited{0.19}。可以计算两个不同的下降量：相对首次适配后表现，下降\calcnum{0.19}；相对初始基座，下降\calcnum{0.14}。前者描述已经获得的局部收益后来保留了多少，后者描述最终是否低于原起点。这两个参照各有意义，不能在不同方法间悄悄交换。

在作者的八组摘要中，首位材料第 1 步和第 8 步的均值分别为\cited{0.3825}、\cited{0.1894}，相减约为\calcnum{0.1931}。同一列的误差还取决于组间差异怎样进入计算。论文附录 B.6 表 5 在第 1 步列出\cited{0.0273}。归档程序对八个组内样本标准差取算术平均，所得数再除以 $\sqrt{64}$，与表中这一格相符。若把八组看作 \cited{64} 条序列的一个样本，用各组均值和方差合并计算，得到的单格标准误约为\calcnum{0.0287}\cite{seal-continual-archive-2025}。两种算法对组间差异的处理不同。计算同一序列从第 1 步到第 8 步的配对差标准误，还需要跨步协方差；这八份摘要只给出逐格统计量。

同样不能从一条下降列推出每次更新都会伤害所有旧任务。其他列有不同起点和变化，系统还可能在某些旧任务上受益。例如，首位材料的平均成绩从第 7 步的\cited{0.1858}回升至第 8 步的\cited{0.1894}。均值总体下降仍可伴随局部回升，也可能掩盖某条序列上的严重退化。Figure~6 支持讨论所示连续编辑序列中的保持风险；它本身不裁决 P6 中“任何更新”的全称判断，也没有把编辑幅度作为独立受控变量。

对于自己的连续更新实验，保留一个任务乘版本矩阵通常比只画最终均值更有信息。每个格子绑定任务标识、模型版本、更新前可见材料与评价次数。若任务逐步加入，后期总体平均值的任务集合也在变化，均值上升可能部分来自加入更容易的新题。可以另报固定旧任务集上的保持、新加入任务上的获得，以及包含二者的业务总体表现。

恢复检查也应回到矩阵。恢复某个检查点后，不必要求每个随机输出逐字相同，但需核对模型工件、适配器、提示、工具与评价范围是否一致。若恢复后表现与当时不同，应先检查这些条件和采样波动，不能立即归因于参数没有真正恢复。相反，文件哈希相同也不能替行为评价背书，它只核对了指定工件。

保持能力可以进入训练目标，也可以只进入发布决策。前一种方法在训练时使用旧任务反馈，会改变优化方向和成本；后一种在候选产生后决定是否接受，不保证优化过程自己学会保持。两者能够组合，但实验需说明回归数据被使用了多少次。反复用同一回归集挑选最优检查点后，再把它称为完全未参与选择的最终评价，会重现前几章的选择问题。

\section{从成功轨迹到梯度：等价需要哪些条件}
\label{sec:ch10-gradient}

Self-Challenging Language Model Agents 的式(2)可作为一个清楚起点。以下先讨论一个固定任务。设采样时使用固定旧策略 $\pi_{\mathrm{old}}$，得到轨迹 $\tau_i$ 和终局奖励 $R_i\in\{0,1\}$。本次求导把已采轨迹和冻结验证器给出的奖励视为固定数据，不对采样过程或奖励函数取全导数。暂不加入优势基线、KL惩罚或其他项，以增大奖励为目标，策略梯度的样本估计可写成\footnote{原文\S4.3式(2)以负的奖励加权对数似然为损失并最小化。本节将对应的策略梯度写为上升方向；比较两种写法时应同时核对损失符号、采样分布与归一化。}
\begin{equation}
 \widehat g_{\mathrm{PG}}=
 \frac{1}{n}\sum_{i=1}^{n}R_i
 \nabla_\theta\log\pi_\theta(\tau_i)
 \big|_{\theta=\theta_{\mathrm{old}}}.
 \label{eq:ch10-pg}
\end{equation}
轨迹对数似然包含相应动作的条件对数概率。奖励为零的轨迹在这个无基线估计式中不贡献梯度，正奖励轨迹贡献增加其似然的方向。这正是原文联系成功轨迹微调的理由\cite{self-challenging-2506}。多任务时还须固定任务分布或权重；若各任务预先分配的采样数不同，池化平均会按这些数量隐式加权，需要核对它与目标权重是否相符。

但“只留下成功轨迹再求平均”通常使用成功数 $m$ 作为分母。若每条轨迹的损失权重一致，则在同一个参数点上
\begin{equation}
 \widehat g_{\mathrm{PG}}=
 \frac{m}{n}\widehat g_{\mathrm{success}},\qquad m>0.
 \label{eq:ch10-normalization}
\end{equation}
方向相同，尺度可能不同。若要让一次普通梯度更新相同，需要同时处理学习率与归一化；若采用自适应优化、正则项、按长度平均或梯度裁剪，单纯把学习率乘一个数也未必完整匹配。若 $m=0$，成功子集平均没有定义，不能把它当作一次正常的微调批次。

看一个只有成功动作与失败动作的教学模型。设成功概率为\calcnum{0.5}，观察四条固定轨迹，其中一条成功、三条失败。对成功动作的对数似然，关于其 logit 的导数为\calcnum{0.5}。无基线、按全部四条平均的奖励梯度为\calcnum{0.125}，仅对成功条目平均为\calcnum{0.5}。两者相差\calcnum{4}倍，这来自分母而非能力差异。此算例是确定的代数检查，没有训练任何模型。

加上一个非零基线之后，失败轨迹也可能贡献样本梯度。仍在同一教学批次中，若基线为\calcnum{0.25}，成功项的权重为\calcnum{0.75}，失败项的权重为\calcnum{-0.25}；按四条平均得到\calcnum{0.1875}。合法基线可以保持期望策略梯度不变，却会改变有限批次估计，这与“每批只训练成功样本”不同。这里展示的是条件变化，不推荐用这四条轨迹选择最优基线。

数据来源也是条件。由当前策略生成的轨迹，和由更强教师生成的轨迹，分布未必相同。直接模仿教师是一个监督学习目标，不能不加修正就叫作当前策略的同策略梯度。原论文式(3)另行讨论强教师轨迹蒸馏，并指出在其初步实验中，包含错误反思的失败轨迹也可能对较弱学生有用\cite{self-challenging-2506}。这说明“失败一律丢弃”不是所有数据训练任务共同遵守的规则。

多轮更新还会改变采样分布。固定成功轨迹训练多轮，是反复提高这份旧数据的似然；每轮重新采样再求奖励梯度，会产生随策略改变的新数据。即使第一步方向一致，后续过程也可能不同。讨论两种算法时，应分别核对生成策略、采样频率、筛选方式、损失权重、正则项和优化器，而不是仅看它们是否都使用二值奖励。

这一推导也帮助理解验证器的地位。若验证器把某种错误轨迹误判为成功，成功筛选与相应正奖励更新都会强化它。改用另一个算法名字不能消除同一反馈错误。反过来，一个可靠验证器提供的稀疏信号也可能让大多数样本权重为零，此时探索和任务难度分配成为问题。第~\ref{ch11}章讨论的任务生成与课程，正是与这里的数据覆盖相接的另一层设计。

\section{偏好、奖励与它们各自的反馈要求}
\label{sec:ch10-reward}

\begin{sloppypar}
不同参数学习目标需要不同材料：监督微调需要目标序列，策略优化需要奖励或优势估计，偏好学习需要比较信息。外部更新同样依赖可靠反馈，不能仅按写入位置确定谁更挑来源。
结合第 4 章对反馈来源的区分，可以讨论以下信号；它们不是适用强度的固定排序：
\end{sloppypar}

\begin{itemize}
  \item \textbf{形式验证器与执行反馈}。执行器在预定任务中可提供可复查的结果信号，Absolute Zero 使用可执行任务反馈开展自博弈\cite{absolute-zero-2505}。执行通过只覆盖相应检查；任务可行性、测试覆盖和真实目标一致性仍需分别核对，执行反馈也不等于形式证明。
  \item \textbf{学习式评审}。能用，但要当成被治理的基础设施来用：
        生成与接受共享未经核对的判断可能形成自证；该综述建议限制评审器修改并保留审计，但这属于治理建议，不能当成已证明充分的技术条件（\cite{self-improvements-survey-2607} §9.1）。
  \item \textbf{内在信号}。多数投票（自一致性）可以当无标签奖励用，
        第 11 章会用一组数字说明它的代价：它只说明\textbf{模型自己同意自己}，
        不说明环境目标完成了。
  \item \textbf{人类偏好}。偏好对提供相对选择信息，与事实正确标签不同。第~\ref{sec:ch10-preferencepair}节据 DPO 原论文解释其目标；本书没有该路线实测，不给跨方法强度排名。
\end{itemize}

\begin{pitfallbox}
\textbf{反例：以为\textbf{会自我编辑}就等于\textbf{不再需要外部任务}。}
原作者指出，主实验的实际适配奖励要求上下文配有下游任务，例如材料附带参考问答，才能评价编辑后的效果（\cite{seal-2506} §5，Context-dependent evaluation）。这限制了该实现直接用于无标注语料。
但不能将其写成所有自我编辑算法的硬约束：本章已介绍的代理奖励改用文本评分，作者还讨论了自生成评价题的方向；改变反馈方式后仍需核对它与真实目标的关系。
换算到工程上：拿到一段无标注文本就去微调，这一步至少还欠一个问题的答案——
\textbf{凭什么说它学会了}。本书的答案仍是第 3 章那套封存集与回归集，
连系统自己改写出来的题目也不豁免。
\end{pitfallbox}

% ---------- 5. 边界 ----------
\subsection{代理奖励省下的是什么}
\label{sec:ch10-proxyreward}

SEAL 的主要外层奖励要实际适配模型并评价，因此成本较高。Appendix~B.10 还考察了一个更便宜的代理：让 GPT-4.1 按若干评分要求评价生成的编辑文本，再用评分作为外层学习信号。Table~9 报告主方法训练约\cited{6}小时、代理版本约\cited{5}分钟，对应的更新后成绩为\cited{47.0}与\cited{45.6}\cite{seal-2506}。这是该论文的特定补充实验。表9训练时间与B.5逐轮计算说明分别保留，不据此重建整个训练的统一总账，也不推定任意任务都能得到同样加速。

两种奖励测量的对象不同。实际适配奖励直接观察某次编辑在给定内层训练与问答中的效果；文本代理先判断编辑看起来是否满足若干标准，希望这种判断与后续学习价值相关。代理节省的是反复进入内层适配的成本，也引入了奖励与真实目标偏离的可能。它没有取消最终任务评价，不能用代理高分本身证明学会新知识。

这项补充结果还说明本章不应将主要实现的一次评价代价写成所有自我编辑方法不可改变的常数。设计者可以改变奖励层级、选择频率和候选筛选方式，但改变之后需重新比较效果、费用和误选率。更快的代理如果倾向选择更长却不准确的材料，可能把节省的搜索时间换成后续训练与回归问题。

合理的使用方式是给代理一个清楚的角色。例如，让它在大量候选中做初筛，再对保留候选执行实际适配；或者用一部分新材料持续校核代理分数与真实收益的关系。这些是可研究的设计，不是本书已实现的机制。初筛漏掉的好候选和放过的坏候选都要进入评价，而非只展示最终选中的案例。

因此，从外部经验进入参数时，验证也可能形成多层结构：来源核对、生成材料检查、廉价代理筛选、真实适配评价、发布回归。层数增加不自动提升可靠性，每一层都可能有成本与偏差。需要说明哪一层负责排除哪一种错误，以及后续评价是否足以发现前面留下的问题。这样的说明比给整个过程贴上“自我监督”或“自我奖励”的名称更有判定价值。


% ---------- 4. 反馈要求 ----------
\subsection{偏好对告诉模型的是相对选择}
\label{sec:ch10-preferencepair}

DPO 原论文从带 KL 惩罚的奖励最大化与 Bradley--Terry 偏好模型出发，把奖励差改写成策略与参考策略的对数概率比，从而直接用偏好对优化策略，不必另训一个显式奖励模型\cite{dpo-2305}。对输入 $x$、被偏好的回答 $y_w$ 和另一回答 $y_l$，其核心损失可简记为
\begin{align}
 z_\theta&=\beta\left[\log\frac{\pi_\theta(y_w\mid x)}{\pi_{\rm ref}(y_w\mid x)}-\log\frac{\pi_\theta(y_l\mid x)}{\pi_{\rm ref}(y_l\mid x)}\right],\\
 \mathcal L_{\rm DPO}&=-\mathbb E\log\sigma(z_\theta).
\end{align}
其中$\beta>0$为KL惩罚系数，$\sigma$为sigmoid函数，期望取自给定偏好对数据；相应策略在观测回答上的概率须使对数比有定义。参考策略参与相对尺度；该机制不自带事实真实性证明。以上是原文第3、4节及式(7)的机制摘要，本书不复述其成绩排名，也没有运行 DPO。


\subsubsection{同一个字母，先看它乘在哪里}

这里的 $\beta$ 对应 DPO 推导中乘在 KL 惩罚前的系数。对固定输入，若参考分布在所考虑的回答上概率为正，且给定奖励使归一化常数有限，原文的理想最优策略为
\begin{equation}
 \pi^*(y\mid x)\propto\pi_{\rm ref}(y\mid x)
 \exp\!\left(\frac{r(x,y)}{\beta}\right),\qquad\beta>0.
 \label{eq:ch10-kl-scale}
\end{equation}
这是奖励优化问题的解，原文再将它代入偏好模型得到 DPO 损失\cite{dpo-2305}。

设一个教学问题只有两个回答，参考概率各为\calcnum{0.5}，奖励差固定为\calcnum{1}。$\beta=\calcnum{1}$ 时，高奖励回答的最优概率约为\calcnum{0.7311}；$\beta=\calcnum{0.5}$ 时约为\calcnum{0.8808}。减小惩罚系数，在这个给定奖励问题中允许更大偏移。这些是理想目标的计算值，有限数据上的 DPO 训练还受优化器、模型表达与停止点影响。

附录~\ref{sec:appc-jitrl} 的 JitRL 则把惩罚系数写成 $1/\beta$，相应偏置含 $\exp(\beta\widehat A)$\cite{jitrl-2601}。两个方法使用同一字母，增大它对相应理想目标的约束强度却有不同方向。比较超参数时，应先核对目标函数中的定义，再比较数值大小。

把这个目标放回开发助手场景，必须先问“偏好”表示什么。两段修复都通过已有测试时，评审可能更喜欢结构清楚的代码；一段通过、一段失败时，偏好可能主要表达正确性；两段都失败时，仍可偏好更接近要求的那段。这三种数据能长得一样，都是一个输入加一对回答，但它们对能力改善的含义不同。只保留胜者和败者标记，会丢掉解释这一差别所需的信息。

本书因此建议在经验记录中保留比较理由、适用要求和已知检查结果。这是材料设计建议，不是 DPO 算法的额外原始定理。如果一条偏好主要依据简洁性，就不要在后续汇总时把它计入“事实正确标签数”；如果依据执行反馈，也要记录执行覆盖了什么。相对选择可以表达有价值的任务需求，但不能自动补上未做的检查。

偏好来源同样影响解释。人类评审、独立模型评审、执行规则及多数投票都能生成一对顺序，顺序的误差结构却不同。让生成候选的同一个模型再选自己喜欢的答案，不一定毫无用途，但缺少外部校核时，只能说明其自有标准趋于一致。将这种过程称为人类偏好学习，或据此宣布事实错误下降，都超出了数据身份。

还有一个容易被名字掩盖的成本问题：不显式训练奖励模型，不等于不用生成候选、取得比较信息和检查更新后表现。若先采样许多回答再选出成对材料，这部分费用发生在优化之前，仍属于经验加工。若复用现成偏好集，可以把本轮新增标注费记为零，但应说明材料来源及与当前任务的差别；不能将历史投入抹去后宣称这条路线天然无反馈成本。

选择监督轨迹还是偏好对，最终应回到自己拥有的证据。可靠示范可以直接说明应输出什么；可靠比较可以说明在两种候选之间偏好哪一个；执行反馈可以说明某个规定检查是否通过。把这些材料转换成训练目标时，可以研究混合使用，但需记录转换规则和损失权重。第~\ref{sec:ch10-gradient}节的一步等价条件也不能直接搬到成对偏好目标上，它们优化的样本表达与目标形式已经不同。


\section{一次参数更新需要保存哪几个版本}
\label{sec:ch10-artifactchain}

参数文件是最终工件，却不是唯一需要保存的东西。一次可复查的更新还包含基座标识、训练材料与来源、数据处理程序、损失定义、优化配置、候选选择规则和评价结果。只保留一个“更新后模型”，很难分辨成绩变化来自哪些材料；只保留训练脚本，也不能保证将来从同名远程模型取得完全相同的基座。

生成材料本身应有版本。假设外层模型将一段来源改写为多条训练样本，后来发现其中一条推断不成立。若没有保存原文到样本的对应关系，就无法确定哪些训练运行使用过它。参数内的影响未必能够按条目精确删除，但训练来源的可追溯性仍然有价值：它至少能支持隔离受影响版本、重新训练候选与明确限制，不能因为难以定位参数就放弃记录材料。

适配器可以改变恢复粒度。若基座固定，更新保存在独立适配器中，系统可以停用该适配器或切回旧适配器；若已经合并进基座并继续训练，恢复就依赖是否保留合并前工件以及后续关系。这里描述的是版本管理条件，不声称停用适配器就能撤销已经发送的工具请求。模型状态恢复与业务补偿仍是两件事。

恢复后继续训练，还可能需要优化器状态、学习率调度与数据进度。仅恢复权重可以得到同一推理起点，却不一定得到相同的后续训练路径。若研究的问题只涉及发布模型的推理行为，未必需要对外分发所有训练中间状态；若声称能够从某个点完整续训，就需要说明恢复范围。把这两种复现目标分开，可以避免既少存必要证据，又无目的地保存大量中间文件。

候选选择过程同样应进入版本链。训练到不同步数可能得到不同候选，最好的一个不是免费出现的。若根据回归集挑选保持最好的检查点，回归集已参与选择；最终保留能力还需在用途清楚的评价范围核对。选择器使用过的每个候选及相应费用不必都进入正文主表，却应能在工件索引中找到。

发布时要核对部署实际加载了什么。模型文件、适配器与提示版本不一致，可以让一个已通过评价的候选在服务中表现不同。若缓存仍保存旧模型生成的计划，或路由将部分任务送到另一版本，总体指标也不再对应单一候选。发布记录应明确版本边界，特别是更新前已开始、更新后才结束的任务归属。

这些记录不保证参数更新有益，它们保证发生问题时能够提出可以检验的问题：是材料有误、选择过拟合、部署版本错配，还是在一致条件下出现了保持下降。把每次下降都叫作灾难性遗忘，会掩盖装置错误；把所有问题都归给装置，又会逃避真实学习失败。版本链的作用正是让两类解释有机会被区分。

\subsection{参数更新后必须通过的三道验证}
\label{sec:ch10-gate}

\begin{sloppypar}
外部持久物与参数都需要回归、发布和风险检查，差别主要在检查对象与可恢复状态。删除记忆不自动撤销派生影响，恢复参数检查点也不自动恢复外部工具状态。以下三项是本书采用的评价框架，不是只对参数方案有效的绝对关卡。
\end{sloppypar}

\begin{enumerate}
  \item \textbf{范围明确的回归评价。}在使用结果选择更新之前，固定需保留的任务、指标与可容忍变化；它只能支持被覆盖范围，不能证明整个模型没有改变。SEAL Figure~6 沿同一编辑序列回测旧任务，并报告逐格标准误，为检查保持变化提供了具体参照。
  \item \textbf{发布决策。}根据新任务收益、旧任务保持、资源和权限约束选择拒绝、受限试用或发布。一个总体均分上升可能掩盖关键子集下降；外部记忆与技能也适用这一原则。
  \item \textbf{与改动相关的风险重测。}综述强调参数更新可能改变既有保证，同时讨论脚手架修改形成持久漏洞的风险\cite{self-improvements-survey-2607}。应检查原保证依赖哪些条件，不声称不动参数就自动安全。
\end{enumerate}

\begin{sloppypar}
E08 的研究问题是同一来源经验进入外部持久物或监督更新后，新任务收益、旧任务保持和生命周期成本如何不同。当前没有可执行的冻结设计与训练资源，不启动参数试验。旧蓝图把“差异区间含零”写成参数保留能力无劣势，这不成立：区间跨零可能同时容纳明显受损与改善。若研究不劣，应事先定义容忍界及相应置信界或检验判据；若研究更新必然退化的普遍主张，则需要针对该主张的有效反例，而非未显著结果。
\end{sloppypar}

\section{同源材料怎样比较外部记忆与参数巩固}
\label{sec:ch10-samesource}

设想一个开发助手已经积累了一批经过复查的修复记录。每份记录含原始需求、错误现象、修改前后的代码、执行反馈与适用条件。团队想知道：下一批任务到来时，把记录放进检索库更合算，还是用它们训练一个适配器更合算？这个问题看似只差写入位置，实际还差材料加工、访问方式、选择预算和部署条件。本节给出的是 E08 的设计草案，不是已经完成的实验。

第一步应冻结共同的来源包，而不是分别给两条路线找最顺手的数据。来源包要说明哪些内容可用，哪些仍未核实，哪些包含未来评价的答案。假如外部路线只有一句摘要，参数路线却得到完整成功轨迹和额外教师解释，两者比较的是整个处理方案，不是单独比较存储位置。反过来，若检索路线可以随时读取原始库，而训练路线只能看到一次截断摘要，也不能据此宣布外部存储保留信息更好。

共同来源不要求输入字节完全一样。检索器需要索引，训练器需要样本，两种接口本来就不同。可比性的关键是记录从来源包到各自工件的变换：删去了哪些字段，生成了什么新内容，是否调用额外模型，是否使用任务答案筛选。把共同来源称为“信息起点相同”比较准确；把两种最后输入称为“信息完全相同”则需要更强证明。摘要可能丢失条件，训练样本可能复制同一事实多次，这些都应作为设计的一部分报告。

一个实用的比较可以分两层。第一层冻结材料加工规则，研究一种具体检索实现与一种具体参数适配实现的表现。这一层容易解释，但结论受所选实现约束。第二层允许各路线在相同开发资源额度下选择自己的加工与配置，比较达到的效果与总费用。这一层更接近实际决策，却把搜索能力也纳入了处理差异。二者都合理，不能用第二层的结果声称已经隔离了写入位置的纯效应。

还应保留一个不吸收新经验的起点。没有它，外部与参数两路即使一样差，也可能被写成两路同样有效。若希望知道收益来自“能访问材料”还是“加工经验”，可再设置直接提供原文的上下文条件；若希望考察两路互补，可加入适配后仍检索的组合条件。增加条件会增加费用和比较数量，应按核心问题选择，不必把所有想法塞进第一次试验。这些条件仍属于比较设计，尚非已有结果。

同源设计还要约束执行器。工具权限、最大交互步数、任务结束条件与输出检查，应尽量一致。如果参数路线使用了更大的上下文或更多重试，它的提高可能来自额外执行资源。反过来，给检索路线增加检索调用却从执行预算中扣除，和另给独立检索额度，是两种不同政策。应在运行前明确比较“相同总费用下谁更好”，还是“相同执行上限下额外机制带来什么”。把每回合上限相同写成总资源相同，会再次漏掉训练和检索的前置投入。

\subsection{划分的不只是题目，还有材料与选择机会}

对来源材料做随机行切分，未必得到独立评价。一个错误模式的原始记录、改写摘要、生成练习与最终测试题，可能共享同一个根案例。若它们散落在训练和评价两侧，模型未见过测试题字面，并不代表没有见过其关键答案。本书建议先按来源案例或任务家族建立归属，再在这个层级上划分；实际单位应服从研究所声称的迁移范围。

可以区分三个用途。材料区用于构造记忆和训练样本；开发区用于选择检索参数、学习率、训练步数和发布门槛；最终评价区用于报告选择完成后的表现。开发区还可包含旧能力回归任务，但一旦根据其结果选择候选，它就是选择数据。重复对最终区查看结果、修改配置、再报告最好成绩，相当于把最终区变成了新的开发区，不能靠保留原文件名维持封存身份。

比较还要说明任务发生的时间。静态比较问的是一份固定经验包在后续任务上的作用；在线比较问的是任务不断完成后，两路分别怎样更新。在线条件下，较强路线可能先解决更多任务，从而得到更多新经验。后期差异包含这种反馈循环，不能都解释为相同材料的写入效果。如果要隔离写入位置，可以先给两路投喂同一条冻结材料序列；如果要研究闭环整体表现，就应允许轨迹分化并完整记录各路新获得的材料。

同一条序列也不自动产生相同曝光量。训练重复读取样本，检索可能只返回少数条目；两路都见过来源包，与每条经验被使用同样多次不是一回事。可同时记录样本训练次数、每次检索返回标识和进入上下文的内容。曝光量本身可能是机制的一部分，不必机械强行配平，但应让读者知道比较中哪些量被固定，哪些量被允许变化。

\subsection{四种结果把失败定位到不同位置}

评价最好至少保留新任务获得、旧任务保持、材料冲突处理与来源追溯四个视角。新任务成功而旧任务下降，提示净收益需要权衡；两者都不变，则可能是材料无用、更新未生效或评价不敏感。不能只凭一张总均分表把这几种情形合并。

材料冲突可以采用一个明确标为教学设计的更新案例：旧规则允许某个操作，新规则在指定版本后取消它。评价不仅问最新条件下是否遵守新规则，也问历史条件下是否仍能解释旧规则的适用范围。若题目只问当前答案，直接覆盖可能很好；若业务要求回溯历史，保留版本条件才重要。这不是偏袒哪种写入位置，而是先把真实任务说明白。

来源追溯也应按实际输出判断。检索到了原文，不代表答案正确引用了它；模型没有直接读取原文，也不代表系统不能通过独立索引给出对应材料。可以分别检查所引来源是否存在、是否支持当前结论，以及是否属于允许使用的版本。把“存在一个引用字符串”当成追溯成功，会让两路都通过一个过弱的测试。

若参数路线失败，至少需要区分适配未能学入、已经学入但后续遗忘、问题表述变化导致调用失败，以及工具执行环节失败。若外部路线失败，则可检查条目是否入库、是否被召回、是否进入上下文、是否被正确应用。诊断粒度不同不意味着只能粗略比较：最终任务指标保持一致，内部诊断按各自机制记录，才有可能解释下一轮该修哪一环。

\subsection{把比较结果变成一次有条件的选择}

考虑同一个修复助手的两类经验。第一类是经常变化的接口说明，回答时还必须给出版本来源；第二类是反复出现的调试步骤，任务文本不同却共享稳定的操作要求。它们可以来自同一个日志库，但业务要求不同。把二者混在一张平均成绩表里选一个永久写入位置，容易掩盖真正的取舍。

对接口说明，设计者可以先检查变化周期与来源要求：更新一条说明之后，旧版本是否仍需可查，撤销时能否定位所有派生材料，服务是否允许临时检索。如果这些条件使外部路线更容易满足当前需求，可以先采用它；这个选择表达的是当前约束下的工程判断，不是参数不能存事实。若以后材料稳定、调用量增长，仍可重新比较适配方案。

对调试步骤，需要检查经验是否真的是可迁移做法。例如“先复现失败，再检查边界输入”可以跨多个问题使用，但一份成功轨迹里也可能夹带偶然路径。训练前可先用外部形式测试它的适用范围，观察哪些任务受益、哪些被误导。这是降低材料不确定性的一种路径，并不保证外部有效的内容写入参数后仍有效。转换成训练样本会改变表达、频率和调用方式，需要新的评价。

假如参数路线在新任务上点估计更高，旧任务的不劣证据不足，费用又可能在较长周期后下降，结论不应被压缩成“参数胜出”。可以将结果写成三个分开的事实：观察到的新任务差值、尚未排除的旧任务损失范围、依赖材料寿命的费用交点。随后根据任务风险与资源条件决定补充评价、限制使用范围或暂不发布。决策可以明确，证据仍可保持不确定。

反过来，外部路线也可能因检索漏召回、上下文膨胀和旧条目干扰而不满足要求。此时应检查失败集中在哪些材料和任务上，能否通过索引、条目更新或适用条件改善。直接从“当前检索实现失败”跳到“必须改参数”，漏掉了实现质量这个中间解释；无限修补检索而从不比较另一种路线，也不能证明当前投入值得。给每条路线限定开发预算，有助于让探索成为可结束的过程。

组合方案更需要明确分工。例如让适配器承担稳定的输出习惯，让外部库提供有版本的事实，表面上兼顾两边优点；但若两处内容冲突，执行者仍要决定听谁。评价中应加入冲突条件，并检查更新一侧时另一侧是否产生旧偏置。两个单独通过的部件拼在一起，不能省去组合后的回归，这与第9章工作流组合的原则一致。

最终的选择记录可以只保留一页，但至少回答四个问题：本次要改善哪类任务；依据哪些材料和评价；哪些费用与风险尚未核实；出现什么观测就重新比较。这样的记录并不要求设计者等所有未知都消失才行动，而是使受限发布的范围与证据相称。本书 E08 仍处于设计讨论阶段，因此这里提供决策方法，不预写一个尚不存在的胜者。后续运行还须保持同样的记录边界。


% ---------- 7. 条件变化 ----------
\section{训练、检索与恢复怎样进入同一份账}
\label{sec:ch10-lifecyclecost}

本节继续使用假想的同源比较，所有金额单位都只是教学记账单位。对路线 $j$，在一次材料更新至下一次材料失效之间，可写成
\begin{equation}
 C_j(N)=A_j+T_j+S_j+V_j+\sum_{i=1}^{N}c_{j,i}+R_j.
 \label{eq:ch10-lifecyclecost}
\end{equation}
$A_j$ 是材料准备，$T_j$ 是训练或索引构建，$S_j$ 是候选搜索与选择，$V_j$ 是发布前回归，$c_{j,i}$ 是第 $i$ 次服务费用，$R_j$ 是该周期内实际发生的恢复与清理费用。名称不重要，边界不能重叠。例如，一次训练运行已经含验证调用时，应将其拆开记账，或在总项内说明所含内容，不能再把同一验证费加一次。

共同材料准备可以有两种报告方式。项目实际总支出只记一次共同投入；比较每条路线独立部署的预算时，各自成本可能都需要包含这份投入。两张账回答不同问题。若同时展示，应显式说明共享分摊，否则读者会看到两个路线之和比项目总账大，却不知道这是重复计费还是两个决策情境。未获账单的失败请求也不能按零费用处理，应保留未知项和可以核实的请求数量。

设教学例中，外部路线前置费用为\calcnum{20}，单次服务费为\calcnum{3}；参数路线前置费用为\calcnum{140}，单次服务费为\calcnum{1}。暂假设质量满足同一要求，周期内无额外更新或恢复，且单次费用恒定。两路总费相等的次数是
\begin{equation}
 N^*=\frac{140-20}{3-1}=\calcnum{60}.
\end{equation}
这只是所列假设下的代数交点，不是本书测出的参数回本点。服务\calcnum{40}次时，两路分别花\calcnum{140}与\calcnum{180}；服务\calcnum{100}次时，分别为\calcnum{320}与\calcnum{240}。若材料只在四十次调用内有效，预计一年有一百次任务也不能据此宣称本次更新能回本，因为后六十次可能已经需要重新训练。

真实系统的单次费用未必恒定。检索索引增长会改变召回与上下文长度，参数服务可能涉及不同硬件、批处理和适配器切换。应先画出累计费用，再检查线性近似是否合适。若参数路线的每次费用并未更低，且前置投入也更高，就不存在由上述摊销机制产生的正交点；它仍可能因质量或其他约束值得选择，但需要另行论证。

费用比较还不能替代质量要求。一条路线很便宜却频繁失败，按请求平均费可能显得优越；只按成功请求费，又可能藏掉失败代价。可同时报告全部尝试总费、完成任务数和预先规定的质量约束。对于必须完成的任务，失败后转交其他系统的费用也属于服务路径。若没有实测转交过程，应把这项留作未知，不用一个方便的固定比例补齐。

恢复费用尤其容易漏掉。外部路线可能需要重建索引、清除派生摘要并重跑受影响任务；参数路线可能需要重新加载旧检查点、清理旧版本生成的缓存并再次回归。恢复工件、恢复服务可用性、确认旧能力恢复，分别消耗不同资源。把文件复制完成的时刻当成事故结束，会低估后两项。相同道理也适用于正常更新中的候选拒绝：没有发布的候选仍然花过训练与评价成本。

时间应与计算费分开记。多张设备并行可以缩短等待，但设备时长之和不会因此自动减少；多个请求的回合时长相加，也不是端到端墙钟时间。面向交付时限，要记录材料可用至候选可发布的实际时间；面向资源预算，则记录调用、令牌、设备时间和已知账单。二者可以一起改善，也可能互相交换。本章未运行 E08，因而不将任何一种交换写成实测优势。

\section{旧能力不劣需要怎样的判据}
\label{sec:ch10-noninferiority}

“没有看到显著下降”与“下降小于可容忍范围”是两个问题。令 $\Delta$ 为新方案减去参照方案在固定旧任务范围上的成功率差，负值表示下降；令 $\delta>0$ 为事先确定的最大可容忍损失。不劣主张关心 $\Delta>-\delta$。若采用置信界决策，需事先确定统计单位、置信水平与计算方法，再检查相应下界是否高于 $-\delta$。不能先看到区间宽度再放大容忍界。

容忍界应来自任务后果，而不是来自能否容易通过。允许普通建议任务有少量变化，与允许关键操作检查漏过错误，可能是完全不同的要求。若某个能力子集不可接受任何实质退化，可以把它列为单独约束，而非在总均分里被其他任务抵消。样本有限时，严格要求可能无法验证；这应如实表现为证据不足，而不是修改术语让结果显得更确定。

下面三个区间仅为教学输入，不是由本书实验估计。设容忍损失为\calcnum{3}个百分点：区间 $[\calcnum{-8},\calcnum{2}]$ 跨过零，但包含超过容忍界的损失，不能据此说不劣；区间 $[\calcnum{-2},\calcnum{1}]$ 的下端高于\calcnum{$-3$}，在相应预定判据下可以支持该范围的不劣，却仍不能说有益；区间 $[\calcnum{1},\calcnum{4}]$ 同时位于零以上，才支持所定义差值方向为正。三者的任务范围、区间构造及多重比较条件仍需一致。

这里的下端可以来自预先规定的双侧区间，也可以来自单侧界；二者对应的尾部错误概率不同，不能只写一个“百分之九十五”而不说明方式。不劣也不等于等效。等效通常需要同时排除两个方向的实质差异，而本节只关心不损失过多旧能力。若新能力增长是另一个必要条件，还应同时检查它；旧能力不劣并不保证更新值得发布。

统计单位应跟随变化来源。对同一份训练材料重复训练若干种子，可以反映训练随机性，但不能替代更多独立材料家族。对一个模型回答许多相似问题，也不能把问题数全当作独立模型更新次数。若目标是新材料家族上的平均保持，需要在相应层级安排覆盖；若只研究一个固定发布版本在当前业务题集的表现，应把结论收窄到这个版本和题集。

配对可以降低某些任务难度差异的影响，但不能消除所有不确定性。两路使用相同题目、相同检查器并记录逐题结果，便于构造差值；相同随机种子在不同模型或不同控制流中未必产生可比随机路径。不能因为文件名里写了同一个种子，就宣布已实现严格的共同随机数实验。实际配对依据应是能够核对的任务和运行条件。

最后，不劣是有范围的评价，不是关于永远不会遗忘的证明。一次发布通过旧任务回归后，后续材料、依赖版本或工具条件变化仍可能改变结果。P6 若讨论更广的更新退化规律，需要明确自己的量词与范围。本节把参数发布问题写成可检验的局部判据，正是为了避免用一个小试替全书的普遍命题作判决。

\section{条件变化：\texorpdfstring{$\theta$}{参数}方案在什么时候才真的占优}
\label{sec:ch10-conditions}

\begin{itemize}
  \item \textbf{同一类做法重复到足够多次}。少量低频技巧可以先比较外部存储；
        若重复使用能摊薄训练与回归投入，参数方案在费用上可能更有吸引力；任务质量、模型能力和部署约束也会影响选择。
        一般摊销思路见第 8、9、13 章，\textbf{但本书没有参数方案的成本实测，所以算不出这个点}。
  \item \textbf{希望学习适配方案的生成方式}。在上述筛选任务范围内，SEAL 的外层策略学会了生成更有效的数据增强与训练配置。
        这为学习适配方案提供了具体实例，迁移范围仍需另行评价。
  \item \textbf{外部持久物已经明显不够用}。外部记忆可能增加检索和上下文，参数方案可能减少这些支出，但仍有训练、服务、路由或检索费用；
        这个差额要先能测出来，才谈得上讨论。
  \item \textbf{反过来，以下情况应先核对额外条件}：更新频繁且需要说清\textbf{为什么变成这样}；
        验证器本身还没冻结；需要按条目追责（第 14 章的归属要求）。
\end{itemize}

% ---------- 8. 能说与不能说 ----------
\section{选择写入位置}
\label{sec:ch10-claims}

选择写入位置需要同时比较来源追溯、更新频率、推理开销与恢复方式。SEAL提供了两种具体路径：将材料改写为训练文本，以及为少样本适配生成配置；连续编辑的热图则将获得与保持放在同一条版本序列上观察。这些证据为设计提供参照，本书E08仍处于设计阶段。

一次选择可以明确记录三个条件：新任务收益需要达到什么范围，旧能力最多容忍多大损失，预计在材料有效期内复用多少次。根据这些条件，再比较外部存储、独立适配器或组合方案。恢复指定版本有助于处置问题，但不能替代能力回归与外部状态核对；容易看到文字的外部方案也需要验收。

\falsifier{E08尚未运行，也未冻结确认性设计。若研究旧能力损失不超过预定界，应在观察结果前规定损失方向、容忍界、统计单位与置信界判据。区间包含零不等于无劣势，有限范围判断也不自动裁决P6的普遍命题。}

\openquestion{E08需要可训练模型、硬件、同源材料与完整比较协议。后续评价应同时保留训练、候选选择、回归、服务与恢复费用，才能检验写入位置在什么条件下值得改变。}

\section*{习题}
\begin{enumerate}[label=10-\arabic*.,ref=10-\arabic*]
  \item\label{ex:ch10-1} 表~\ref{tab:ch10-seal} 的少样本设置先筛选 Oracle TTT 可解的任务，因此 Oracle 为 \cited{100}，
        而 SEAL 是 \cited{72.5}。这个上界会怎样改变对 \cited{72.5} 的解读？
        请说明它为何不代表未筛任务上的普遍上界；若要区分编辑生成不足与适配过程受限，各需补什么对照？
  \item\label{ex:ch10-2} 比较表~\ref{tab:ch10-seal} 的单段设置与表~\ref{tab:ch10-cpt} 的两档多段设置。SEAL 相对 GPT-4.1 合成数据的差值依次为 \calcnum{+0.7}、\calcnum{$-1.2$} 和 \calcnum{$-2.8$} 个百分点。请先指出训练方式及评价范围的变化。
        若语料规模可以继续加大，你会在什么位置停止使用自生成数据？
        为了在事先就定好这个位置，需要预注册哪几个量？
  \item\label{ex:ch10-3} 第~\ref{sec:ch10-gate}节的三项检查分别怎样用于外部记忆与独立适配器？请区分恢复文件、恢复行为与撤销外部影响。
  \item\label{ex:ch10-4} \cite{self-challenging-2506} §4.3 给出“二值奖励下一步 REINFORCE 等价于拒绝采样微调”。
        请举一个使这个等价不成立的最小改动（改奖励、改步数，还是改数据筛选），
        并说明改动之后你比较的到底是“算法”还是“目标”。
  \item\label{ex:ch10-5} 教学设定在上述知识注入报告条件下串行评估 \calcnum{100} 个候选自我编辑，每个候选进行一次适配与评价，按每次 \cited{30--45} 秒估算要花多久；
        是否可以直接估算同样时间能写多少条外部经验？列出还缺的硬件、任务、并发与成本范围信息。
\end{enumerate}

\section*{常见误解}
\begin{itemize}
  \item \textbf{“写参数是写持久物的加强版，能早就早”。}
        它改变学习与部署条件，效果和费用需要评价。保留检查点或兼容适配器可以恢复指定状态，不能笼统称不可逆；外部更新也需要风险检查。
  \item \textbf{“自生成数据一定比外部数据好”。}
        三个所列设置中，单段点估计较高；到 $n=200$ 与 $n=2067$ 两档，
        它的点估计低于 GPT-4.1 生成数据条件，不能据此作所有任务上的强弱排序。
  \item \textbf{“用了强化学习就不是监督微调了”。}
        该联系只覆盖明确的采样、加权与一步更新条件；优势基线、KL项、离策略数据或多轮重采样都会改变比较。
  \item \textbf{“改掉一个事实就算学会了这个能力”。}
        事实更新与能力学习是两件不同的事，验收方式也不同：
        前者可用有版本条件的问答评价，后者要明确所声称任务范围并检查未参与学习的任务；是否需要分布外评价取决于迁移主张。
\end{itemize}

\section*{进一步阅读}
\begin{itemize}
  \item \cite{seal-2506}：本章两条主线的出处。Table 1 是“路子本身可学”的证据，
        Table 2 比较不同训练设置；§5 的 Limitations 里有连续更新下带数字的热图、
        “一次自我编辑评价 \cited{30--45} 秒”的成本说明。知识注入逐轮计算、少样本 TTT 计算条件与代理奖励时间分别见 Appendix~B.5、A.4与B.10。
  \item \cite{self-challenging-2506}：§4.3 式 (2)(3) 给出蒸馏设置与自改进设置下的等价与不等价，
        提供特定目标关系的一手表达，不能推成强化学习普遍多余。
  \item \cite{absolute-zero-2505}：§6 与 Figure 34 的 "uh-oh moment"，
        作者自陈未解决这类系统的安全治理——它是本章“安全重测不可省”的动机，不是它的证据。
  \item \cite{self-improvements-survey-2607}：§9.1 的从快探索到慢巩固，
        为写入位置与验证边界提供观点来源；其治理建议不等同于本书已测规律。
  \item \cite{dpo-2305}：奖励目标、最优策略与偏好损失的联系；其 KL 系数定义可与附录~\ref{sec:appc-jitrl} 的 JitRL 对照。
  \item 附录~\ref{app:rl}：策略梯度与信用分配的推导，与本章的一步样本更新相衔接。
\end{itemize}
